% ----------------------------------------------------------
% Conclusão
% ----------------------------------------------------------
\chapter{Conclusão}

O desenvolvimento deste trabalho mostrou como as tecnologias digitais, especialmente os aplicativos móveis capazes de ter um papel estratégico na valorização, no registro e na democratização do patrimônio cultural, principalmente em contextos regionais. Com base em uma fundamentação teórica sólida e na análise de experiências anteriores, ficou evidente que soluções digitais são capazes de ajudar a superar barreiras de acesso, incentivar o envolvimento das comunidades e reforçar identidades locais, especialmente em áreas que costumam ser deixadas de lado pelas políticas culturais mais centralizadas.

O aplicativo criado neste projeto, desenvolvido em Flutter, reúne funções como mapeamento com geolocalização, cadastro colaborativo, acessibilidade e integração de imagens com descrições dos pontos culturais. A arquitetura foi pensada com base em boas práticas e nas referências estudadas, o que garante flexibilidade para futuras melhorias, como a conexão com bancos de dados reais ou até a adição de tecnologias mais interativas, como a realidade aumentada. Ao combinar usabilidade com participação popular e representatividade cultural, a proposta amplia o alcance das manifestações culturais e facilita o acesso à informação para moradores e visitantes.

A análise de exemplos similares, como o Vem Para Estância e plataformas públicas como o Mapa da Cultura, ajudou a reforçar a importância de sistemas colaborativos, fáceis de usar e adaptados às realidades locais. Este projeto se aproxima dessas iniciativas ao colocar a participação dos usuários como parte essencial do processo de construção e atualização dos dados, contribuindo para uma cultura digital mais inclusiva, viva e em constante transformação.

Em resumo, pode-se concluir que o uso de aplicativos móveis voltados ao mapeamento cultural é um caminho real e necessário para fortalecer o direito à memória, valorizar a diversidade e aproximar as novas gerações de suas raízes culturais. Pesquisas futuras poderão explorar mais profundamente os impactos dessas ferramentas e testar novos recursos de acessibilidade e interação, consolidando a tecnologia como uma parceira essencial das políticas públicas de cultura no Brasil.
